\begin{pastexam}[label={px:ntut-107-prog-03}]{北科大資工所 107 學年度 第 3 題}
（程式填空）

\textbf{Problem 3 [30\%, each 3\%]}

The following program can transfer an infix expression into a postfix expression according to the
BODMAS (Brackets, Orders, Division, Multiplication, Addition, and Subtraction) rule that is used for
ordering mathematical operations. In this program, suppose that brackets are not used. The
operators include "+, $-$, *, /", and the operands consist of "a $\sim$ z". The sample inputs of
the program are: \{\texttt{"a+b*c-d/e\^{}f", "a*b-c/d\^{}e*f+g*h"}\}, and their corresponding
outputs are \{\texttt{"abc*+def\^{}/-", "ab*cde\^{}/f*-gh*+"}\}. Please trace this program and fill
the 3-1$\sim$3-10 blanks with correct answers.

\begin{minted}[fontsize=\small]{c}
#include <stdio.h>
#define SIZE 20
typedef enum {false=0, true =1} bool;
bool isEmpty(char stack[], int top) {
    if (__problem 3-1__ <=-1) return true;
    else return false;
}
bool push(char stack[], int *topp, char element) {
    stack[__problem 3-2__(*topp)] = element;
    return true;
}
bool pop(char stack[], int *topp, char *element) {
    if (isEmpty(stack, (*topp))) return false;
    (*element) = stack[(*topp)__problem 3-3__];
    return true;
}
void popAll(char stack[], int *topp, char postfix[],
            int *indexp) {
    while (!__problem 3-4__(stack, (*topp))) {
        postfix[(__problem 3-5__)++]
                            =stack[(*topp)--];
        printf("%c", postfix[(*indexp)-1]);
    }
}
bool isOperand(char c) {
    if ((c>='a')__problem 3-6__(c<='z')) return true;
    else return false;
}
int opPriority(char c) {
    if (c=='^') return 3;
    else if (c=='*' || c=='/') return 2;
    else if (c=='+' || c=='-') return 1;
    else return 0;
}
bool shouldPopOp(char x, char y) {
   if (opPriority(x)__problem 3-7__opPriority(y))
       return true;
   else return false;
}
void operatorProcess(char stack[], int *topp,
                 char postfix[], int *index, char c) {
    char elem;
    bool nonEmpty;
    nonEmpty = pop(stack, topp, &elem);
    while (nonEmpty&&shouldPopOp(elem, c)) {
        postfix[(*index)++] = elem;
        printf("%c", __problem 3-8__);
        nonEmpty = pop(stack, topp, &elem);
    }
    if (__problem 3-9__) push(stack, topp, elem);
    push(stack, topp, c);
}
void operandProcess(char postfix[], int *indexp,
                    char c) {
    postfix[(*indexp)++] =__problem 3-10__;
    printf("%c", c);
}
void infixToPostfix(char infix[], char postfix[],
                    int size) {
  int i=0, index=0, top=-1;
  char stack[SIZE];
  for (i=0; i<size; i++) {
    if (isOperand(infix[i])) {
      operandProcess(postfix, &index, infix[i]);
    }
    else {
     operatorProcess(stack, &top, postfix,
                     &index, infix[i]);
    }
  }
  popAll(stack, &top, postfix, &index);
}
void test01() {
    char infix1[] = "a+b*c-d/e^f";
    char infix2[] = "a*b-c/d^e*f+g*h";
    char postfix[30];
    int size = 15;
    infixToPostfix(infix1, postfix, size);
    infixToPostfix(infix2, postfix, size);
}
\end{minted}
\end{pastexam}

\begin{pastsoln}
答案：
\begin{center}
\begin{tabular}{*{5}{l}}
\toprule
3-1 \texttt{top} & 3-2 \texttt{++} & 3-3 \texttt{--} & 3-4 \texttt{isEmpty} & 3-5 \texttt{*indexp} \\
3-6 \texttt{\&\&} & 3-7 \texttt{>=} & 3-8 \texttt{elem} & 3-9 \texttt{nonEmpty} & 3-10 \texttt{c} \\
\bottomrule
\end{tabular}
\end{center}

\textbf{原題的兩個問題}。 第一，題目說運算子只有 $+$、$-$、$*$、$/$，但範例和 \texttt{opPriority} 都有
\verb|^|（次方，優先序 $3$），照程式和範例把 \verb|^| 當成運算子。 第二，\texttt{test01} 對兩個字串都傳
\texttt{size = 15}，但 \texttt{infix1} 只有 $11$ 個字元（陣列長度 $12$，含結尾的 \verb|'\0'|），迴圈會讀到
\verb|infix1[11]|（\verb|'\0'|）和陣列外面的 \verb|infix1[12..14]|，是未定義行為；讀到 \verb|'\0'| 時它不是
小寫字母，會被當成運算子送進 \texttt{operatorProcess}，把堆疊整個彈光。 要得到題目給的輸出，
\texttt{infix1} 要傳 \texttt{11}（或用 \texttt{strlen}）。 這不影響填空的答案，下面照字串的實際長度追蹤。

\textbf{堆疊的約定}（資結書第~6.1~節的陣列堆疊）：\texttt{top} 是頂端元素的索引，空堆疊時是 $-1$。
\begin{itemize}
\item \textbf{3-1 \texttt{top}}：\texttt{isEmpty} 收到的是 \texttt{top} 本身（不是指標），$\texttt{top} \le -1$ 就是空的。
\item \textbf{3-2 \texttt{++}}：推入要先把 \texttt{top} 加一再存，\texttt{stack[++(*topp)]}。 空堆疊
$\texttt{top} = -1$，第一個元素存在 \texttt{stack[0]}。 寫 \texttt{(*topp)++} 的後置形式放不進這個空格，
而且會先寫到 \texttt{stack[-1]}。
\item \textbf{3-3 \texttt{--}}：彈出是先讀頂端再減一，\texttt{stack[(*topp)--]}，和 \texttt{popAll} 裡的寫法一樣。
\item \textbf{3-4 \texttt{isEmpty}}：\texttt{popAll} 在堆疊不空時一直彈，\texttt{!isEmpty(stack, (*topp))}。
\item \textbf{3-5 \texttt{*indexp}}：\texttt{indexp} 是指向輸出位置的指標，要遞增的是它指到的整數；下一行
\texttt{postfix[(*indexp)-1]} 也證實了這一點。 寫 \texttt{indexp} 會把指標本身加一，型別也不對。
\item \textbf{3-6 \texttt{\&\&}}：小寫字母要同時滿足 $\ge$ \texttt{'a'} 和 $\le$ \texttt{'z'}。
\item \textbf{3-7 \texttt{>=}}：遇到新運算子 \texttt{c} 時，堆疊頂端的運算子 \texttt{x} 優先序\textbf{大於或等於}
\texttt{c} 就先輸出（資結書虛擬碼~6.8）。 等號是必要的：同級的 $-$、$+$ 是左結合，\texttt{a-b+c} 要先算
\texttt{a-b}。 用範例驗證：寫成 \texttt{>} 時，第一個範例會輸出 \texttt{abc*def\^{}/-+}，第二個輸出
\texttt{ab*cde\^{}f*/gh*+-}，都和題目給的不同；寫成 \texttt{>=} 才和題目一致。
\item \textbf{3-8 \texttt{elem}}：剛彈出、剛寫進 \texttt{postfix} 的運算子就是 \texttt{elem}（寫成
\texttt{postfix[(*index)-1]} 也對）。
\item \textbf{3-9 \texttt{nonEmpty}}：迴圈結束有兩種原因。 堆疊彈空了（\texttt{nonEmpty} 是 false），什麼都
沒拿出來；或是彈出的 \texttt{elem} 優先序比 \texttt{c} 低，不該輸出，這時 \texttt{nonEmpty} 是 true，必須把
它推回去。 所以條件是 \texttt{nonEmpty}。
\item \textbf{3-10 \texttt{c}}：運算元直接寫進輸出，下一行印的也是 \texttt{c}。
\end{itemize}

\textbf{追蹤第一個範例} \texttt{a+b*c-d/e\^{}f}（只列運算子）：
\begin{center}
\small
\begin{tabular}{clll}
\toprule
讀入 & 動作 & 堆疊（右邊是頂端） & 到目前的輸出 \\
\midrule
\texttt{+} & 堆疊空，推入 & \texttt{+} & \texttt{a} \\
\texttt{*} & 彈出 \texttt{+}，$1 \ge 2$ 不成立，推回；推入 \texttt{*} & \texttt{+ *} & \texttt{ab} \\
\texttt{-} & 彈出 \texttt{*}（$2 \ge 1$）、\texttt{+}（$1 \ge 1$）都輸出；推入 \texttt{-} & \texttt{-} & \texttt{abc*+} \\
\texttt{/} & 彈出 \texttt{-}，$1 \ge 2$ 不成立，推回；推入 \texttt{/} & \texttt{- /} & \texttt{abc*+d} \\
\verb|^| & 彈出 \texttt{/}，$2 \ge 3$ 不成立，推回；推入 \verb|^| & \verb|- / ^| & \texttt{abc*+de} \\
結束 & \texttt{popAll} 依序彈出 \verb|^|、\texttt{/}、\texttt{-} & & \texttt{abc*+def\^{}/-} \\
\bottomrule
\end{tabular}
\end{center}
和題目給的輸出相同。 第二個範例同樣得到 \texttt{ab*cde\^{}/f*-gh*+}。 附帶一提，用 \texttt{>=} 時
\verb|^| 也被當成左結合（\verb|a^b^c| 會輸出 \verb|ab^c^|），數學上次方通常是右結合；範例裡沒有連續
的 \verb|^|，所以不影響本題。
\end{pastsoln}
